\section{Gemini 4 Carbon、サイバー分野と長期タスクを訴求}
Google DeepMindは、次のフロンティアモデルとしてGemini 4 Carbonを発表した。脆弱性発見や修正を含むサイバーセキュリティ、長い手順を要するタスクで最先端の性能を主張する。Grok収集によれば、Artificial Analysis Intelligence Index v5.0で59点、SEC-bench Proで88.2\%を記録したとされる。初期提供は防御側を対象とするFairwind Programに限定し、安全対策を改善しながら拡大する方針。導入価格は100万トークンあたり入力2米ドル、出力10米ドルと説明された。数値と性能評価は発表資料に基づき、独立した再検証は本号では行っていない。
\dailyxsource{Google DeepMind (@GoogleDeepMind)}{https://x.com/GoogleDeepMind/status/2115489356437651654}

\section{Genie 4の世界モデルを巡る議論が継続}
前日10月27日に発表されたGoogle DeepMindのGenie 4について、観測期間中もデモや活用先を巡る投稿が続いた。リアルタイムに操作できる世界シミュレーションとして、音声生成、ステアリング、Geminiとの統合が特徴として挙げられる。ただし、公式発表そのものは今回の観測窓外であり、窓内の反応と混同できない。Grokが例示した関連X投稿は10月29日07時01分頃と記録されており、厳密には今回の終了時刻を約1分越えている。APIやAI Studioへの展開状況は今後の確認事項である。
\dailyxsource{Jason Stephen (@jastephx)、関連投稿（観測窓外約1分）}{https://x.com/jastephx/status/2115564655779381693}

\section{Amazon Novaモデル群の運用縮小・再編報道}
AmazonがNovaの複数モデルを最小限の保守に移し、AI開発組織を再編してフロンティアモデルへ資源を集中するとの報道・投稿が共有された。対象候補にはPremier、Omni、Reel、Canvasなどが挙げられ、ロボティクス等への重点移行も示唆されている。しかし、Grok収集ではAmazonまたはAWSによる公式一次発表を確認できておらず、具体的な製品影響も不明である。紹介されたX投稿は07時06分で、観測窓を約6分越える。現時点では確定した事業方針ではなく、二次報道に基づく主張として扱う。
\dailyxsource{ShushantoKumar (@ShushantoK)、未確認報道（観測窓外約6分）}{https://x.com/ShushantoK/status/2115565773057818808}

\section{DiffuSpace、拡散言語モデルCloverを公開}
DiffuSpaceは、オープンウェイトの拡散言語モデル「Clover-35B-A3B-Preview」を公開したと報告された。総パラメータ約350億、稼働パラメータ約30億のMoE構成で、従来の自己回帰型とは異なる生成方式を採る。単一GPUのBF16実行で毎秒680トークン、SWE-bench Verifiedで73.6という数値も紹介されているが、いずれも提供側の主張であり独立測定値ではない。Apache-2.0ライセンスでの公開が記され、モデルの実測速度、品質、実行条件の比較が今後の論点となる。Grokは直接のPrimary X投稿URLを特定していない。
\dailyxnote{モデル公開先はGrok原文に記載。一次X投稿のURLは未特定のため掲載しない。}

\section{Articul8、業務判断向けBiviaを発表}
Articul8はマルチモーダル意思決定モデル「Bivia」を発表した。個別の質問への応答にとどまらず、文書、画像、音声、動画、センサーデータを扱いながらワークフロー全体を進め、不確実な判断では人間にエスカレーションする設計を訴求する。Claude Codeなどとの連携や、100万入力トークンあたり0.042米ドルという価格も投稿で紹介された。ただし、公式ドキュメントや独立した運用評価は今回の収集では十分確認されていない。具体的なAPI仕様や業務環境での再現性を見極める必要がある。
\dailyxsource{Articul8 AI (@Articul8\_AI)}{https://x.com/Articul8_AI/status/2115558202029007195}

\section{Liquid AIの小型オープンウェイトLFM3}
Liquid AIのLFM3-1.7Bと関連モデルに関する投稿では、端末上での実行や小型エージェントへの適用が話題となった。Grok収集では、4ビット量子化時に約1GBの容量、256Kトークンのコンテキスト長などの特徴が挙げられている。小規模モデルでも長い文脈や実用的な動作環境を目指す動きとして注目される一方、実機での速度、品質、ライセンス条件は利用前の確認が必要である。例示されたX投稿は07時00分と記録されており、観測窓の終了境界に位置する。
\dailyxsource{Marwan Moftah (@MarwanMoftah)}{https://x.com/MarwanMoftah/status/2115564405106630873}

\section{Vangrid、Physical AI向け実世界データ収集を提案}
Vangridに関するコミュニティ投稿では、スマートフォンで取得した実世界のデータをオンチェーンで検証し、報酬付きの依頼を通じてAIエージェント向けのデータを集める構想が紹介された。物理環境の正解データを収集する方法として関心を集めたが、一次公式資料や独立検証は限定的である。投稿中の取得件数などの数値も裏付けが不足している。現段階では仕組みの提案や宣伝を中心とする未検証情報として読み、APIの実在、データ品質、検証方法を今後の確認事項とする。
\dailyxsource{コミュニティ関連投稿（未検証）}{https://x.com/0x_zoda/status/2115276307915313234}

\section{DeepSeek V5のモデルIDを巡る未確認リーク}
DeepSeek V5について、テスト用のHarnessに「deepseek-v5」「deepseek-v5-open」「deepseek-v5-next」などの識別子が現れたとするX投稿が話題になった。投稿者は新モデルの兆候として紹介したが、Grokの収集では公式発表や一次資料は確認されておらず、公開時期や実際の仕様を示す根拠もない。Grokによる否定的な反応も記録されている。現時点でリリースを示す事実として扱わず、公式情報が出るまで未確認のリーク主張にとどめる。
\dailyxsource{Md Ismail Šojal (@0x0SojalSec)、未確認投稿}{https://x.com/0x0SojalSec/status/2115561842588471541}
